\section*{Appendix}
%\subsection{pCUT}
%(\mandal{I see that the results on this one particle dispersion have been removed.  can one formulate it for some very small values of $J_3$ for $J_2=0$? Then also it portends some three-dimensional results and hence can be included.})
%Here, we aim to obtain the equivalence of the current Hamiltonian ~\ref{Hamiltonian} to that of the Toric-Code Model~\cite{vidal2009} in the absence of dipole-dipole interaction, $J_2$ and interlayer coupling $J_3$. We show below that the physics of the current model can be exactly described by the TCM under $K_x\neq 0$, $K_{y,z}=0$. Interestingly, the equivalence can only be established for specific conditions that $J_0=J_1$, rest of the cases, this may not be the scenario. \\
%\indent In the low field limit $J_{0,1} \ggg K$, for 2D intraplane structure in the absence of dipole-dipole interaction, we calculate the one-particle dispersion within the perturbative Continuous Unitary Transformation(pCUT) formalism. Here we shall outline the implementation details of pCUT and give the hopping amplitudes and the corresponding ground state energy up to order 6.
%From the earlier studies we know that the low-field case hosts a deconfined phase with gauge charge excitations coming from the gauge term $J_0$ similar to the excitations in the Toric-code model(TCM) except the fact that in TCM there are two kinds of excitations one is the magnetic charge and the other being the gauge charge corresponding to two types of interactions. It should be noted that the true one-particle excitation on a lattice can be achieved only in the open boundary conditions where one can separate the excitation pairs ideally without any energy cost to infinite distance effectively becoming a one-particle sector.  However, It is to be noted that even in the absence of dipole-dipole interaction the spectrum of unperturbed Hamiltonian isn't equidistant, and to make pCUT applicable we, therefore, choose the parameters so that within the set parameters it isn't the case anymore.\\
%\begin{align}
%e_1&=1 + K (-2 \cos(q_x - q_y) - 2 \cos(q_x + q_y)) + \\ \nonumber
% & K^2 (2 - 2 \cos(2 q_x) - \cos(q_x - q_y) - 2 \cos(2 q_y)- \\ \nonumber
% & \cos(q_x + q_y)) + \\ \nonumber
% & K^3 (-3 \cos(q_x - 3 q_y) - 3 \cos(3 q_x - q_y)  - 3 \cos(3 q_x + q_y) \\ \nonumber
% & - 3 \cos(q_x + 3 q_y)) \\ \nonumber
%    & K^4 (\frac{35}{2} - 2 \cos(2 q_x) - \frac{15}{2} \cos(4 q_x)   - 5 \cos(q_x - 3 q_y) % \\ \nonumber
% &- \frac{25}{4} \cos(2 q_x - 2 q_y) - \cos(q_x - q_y) -  5 \cos(3 q_x - q_y) - \\ %\nonumber
%   & 2 \cos(2 q_y) - \frac{15}{2} \cos(4 q_y) - \cos(q_x + q_y) \\ \nonumber & - 5 \cos(3 %q_x + q_y) -
%    \frac{25}{4} \cos(2 q_x + 2 q_y) - 5 \cos(q_x + 3 q_y))  \\ \nonumber
%\end{align}
%\indent The configurational energy sets for a single given plaquette(four spins) in the absence of magnetic field $K$ are $-J_0 + 2J_1$, $-J_0+2J_1$, $J_0$, Clearly for two cases the spectrum is equidistant, i.e, for $J_1=0$ and $J_1=J_0$ with corresponding $-3J_0$, $J_0$ sets. One may shift the energy to the energies or renormalize it to give the energies $-J_0$ and $J_0$ respectively. Thus we finally get an equidistant unperturbed spectrum. We use a single period cluster lattice to obtain the physical properties such that the order of perturbative expansion, $n$,  doesn't account finite-size effects. In other words, $L \ge n+2$. \\
%\indent The ground state energy per spin $\text{e}_{\text{GS}}$ up to order 6 is given by,
%\begin{eqnarray}
%e_{GS}=-J_0-2J_1 - \frac{1}{2}K^2 - \frac{15}{8}K^4 - \frac{147}{8}K^6
%\end{eqnarray}
%and similarly we give the one particle dispersion $e_{1}$ up to order 4, order 5 and 6 are too lengthy to be given here but will be effectively given when the one particle gap is obtained at the minimum K-point, i.e, $\Gamma$. The $\Delta$ is given by,
\subsection{Real space Linear-Spin wave theory}
\label{RSLSWT}
Following Ref.~\cite{moessner2019} the dynamic structure factor was obtained within the semi-classical approach of Linear-spin wave theory calculated in the real space. Below is a brief description of the method and the implementation details of our code. Since we know the classical ground states obtained from the previous studies We, therefore, perform the Holstein-Primakoff transformation upon these ground states, neglecting terms beyond quadratic in $a,a^{\dagger}$ we diagonalize the Hamiltonian in the basis of $a,a^{\dagger}$ in the real-space. Then using the dynamical matrix method approach the obtained eigenvalues($\omega_\alpha$) and eigenvectors($\psi^{i}_{\alpha}$) are used to calculate the DSF as,
\begin{align}
 S^{xx}(\kappa,\omega)&= \frac{1}{2N}\sum_{i,j=1}^{N}e^{-ik\cdot(r_i - r_j)}\\ \nonumber
 & \times \sum_{\alpha}(\eta_{i}^{x}\cdot\psi_{\alpha}^{i})(\eta_{j}^{x}\cdot\psi_{\alpha}^{j})^{\star}\delta(\omega-\omega_\alpha)\text{sgn}(\omega_\alpha)
\end{align}
where $\eta_{\alpha}^{i}$ are the matrix elements of the rotational matrix obtained after projecting the local $z$-projected classical ground states onto the global coordinate frame. Since at the quadratic level most of the ground states haves same Hamiltonian elements and hence the averaging of the structure factor here is not necessary. The results in sec~\ref{dynamic} are calculated on a 16*16 lattice with periodic boundary conditions invoked.
